
\input texsis
\paper

\overfullrule=0pt
\def\frac#1#2{#1\over #2}
%
\def\bmatrix#1{\left[ \matrix{#1} \right]}
\def\be{{\beta}}
\def\toone{{t+1}}

%\def\frac#1#2{#1\over #2}
%\magnification=1200
\overfullrule=0pt
%


\centerline{\bf Optimal Taxation Models}



$$ E_t \sum_{j=0}^\infty \beta^t L(T_t)  $$

$$ L(T) = .5 T^2. $$


$$ T_t + a_t = R^{-1} a_{t+1} +  G _t   , \quad t \geq 0 $$

or

$$ T_t + a_t(z_t) = \int  q(z_{t+1} | z_t ) a_t(z_{t+1}) d z_{t+1}

where
$q(z_{t+1} | z_t)$ is a pricing kernel for one-period  state contingent claims and $b_t(z_{t+1})$ is   a vector of state-contingent debts due in state $z_{t+1}$
 that are chosen at time $t$.
We exogenously set the pricing kernel to 
$$ q(z_{t+1} | z_t ) = \beta \phi(z_{t+1} \ z_t) ,  $$
where $\phi(z_{t+1} \ z_t)$ is the probability density of $z_{t+1}$ conditional on $z_t$.  
This specification of the pricing kernel
implies that the price of a risk-free one period bond must be
$$ \int q(z_{t+1} | z_t) d z_{t+1} = \beta \int \phi(z_{t+1} \ z_t) = \beta $$
so that the gross interest rate on one-period risk free loans is  $\beta^{-1}$.


$$  T_t =  (1-\beta)
\left[ \sum_{j=0}^\infty \beta^j E_t y_{t+j} - b_t\right]. \EQN sprob9 $$

$ 1 - \beta = {r \over 1+r} $

 The consumer maximizes
\Ep{sprob1} by choosing a consumption, borrowing plan
 $\{c_t, b_{t+1}\}_{t=0}^\infty$ subject to the sequence of budget constraints
$$ T_t + a_t = R^{-1} a_{t+1}  + G_t  $$
where $\{y_t\}$ is an exogenous
 stationary endowment process, $R$ is a constant gross
risk-free interest rate, $a_t$ is one-period risk-free  debt  maturing at
$t$, and $b_0$ is a given initial condition. It is important to note that for $t \geq $, $b_t$ is chosen at time $t-1$ and that $b_0$
is a given initial condition at time $t=0$.  We shall assume
that $R^{-1} = \beta$.  For example, we might assume that the endowment
process has the state-space representation
$$ \EQNalign{ z_{t+1} & = \check A z_t + \check C w_{t+1} \EQN sprob15;a \cr
               y_t & = \check G  z_t \EQN sprob15;b \cr}$$
where $w_{t+1}$ is an i.i.d.~process with mean zero and
identity contemporaneous covariance matrix, $\check A$ is a stable matrix,
its eigenvalues being strictly below unity in modulus, and
$\check G$ is a selection vector that identifies $y$ with a particular
linear combination of the $z_t$.
We impose the following condition on the
consumption, borrowing plan:
$$  % \lim_{T \rightarrow +\infty}
 E_0 \sum_{t=0}^\infty \beta^t b_t^2 < +\infty. \EQN sprob3 $$
This condition suffices to rule out Ponzi schemes.
The {\it state\/} vector confronting the household at $t$  is
$\left[\matrix{b_t & z_t\cr}\right]'$, where $b_t$ is its one-period debt falling
 due at the beginning of period $t$
and $z_t$ contains all variables useful for
forecasting its future endowment.


$$ T_t = {r \over 1+r}
\left[ \sum_{j=0}^\infty \beta^j E_t y_{t+j} - b_t\right]. \EQN sprob9 $$
Equation \Ep{sprob8} or \Ep{sprob9} asserts that tax collections  equal
economic {\it income\/}, namely,  a constant
marginal propensity to consume or interest factor ${r \over 1+r}$ times
the sum of nonfinancial wealth $
\sum_{j=0}^\infty \beta^j E_t y_{t+j}$ and financial
wealth $-b_t$.   Notice that we can use a version of  our formula \Ep{discount1} in  \Ep{sprob8} or \Ep{sprob9} to obtain
$$  T_t  = (1-\beta) [ \check G(I-\beta \check A)^{-1} z_t - b_t ]  ,  \EQN eqn:lucascritiquecons $$
a ``consumption function'' that   represents
$T_t$  as a function of the {\it state\/} $[b_t, z_t]$ that
 confronts the household, where  from \Ep{sprob15} $z_t$ contains present
information that is useful for forecasting government expenditures.

If we subtract equation  \Ep{sprob16;b}  evaluated at time $t$ from
equation  \Ep{sprob16;b}  evaluated at time $t+1$, we obtain
$$ a_{t+1}- a_t = \check G (I-\beta \check A)^{-1} (z_{t+1} - z_t) - {\frac{1}{1-\beta}}(T_{t+1} - _T ) .$$
Substituting $z_{t+1} - z_t = (\check A - I )z_t + \check C w_{t+1}$ and equation \Ep{sprob16;a} into the
above equation and rearranging gives
$$ a_{t+1} - a_t =\check G (I - \beta \check A)^{-1} (\check A - I) z_t . \EQN debt_evolution $$
This equation indicates how $b_{t+1}$ is predetermined at time $t$ as a function of $b_t$ and $z_t$.

\section{Complete Markets}


 To construct a financial plan that supports constant consumption, temporarily take $\bar c$ as known, substitute \Ep{guess101} into
\Ep{stareatit101} and solve for a state-dependent debt function $b(z_t)$:
$$ b(z_t, \bar c) =  \check G(I-\beta \check A)^{-1} z_t - {\frac{1}{1-\beta}} \bar c . \EQN guess102 $$
To pin down the constant $\bar c$, we assume that the consumer enters time $0$ owing debt $\bar b_0$ and that $\bar c$ solves  the following version of \Ep{guess102} at time $t=0$:
$$ \bar c = (1 -\beta) [ \check G(I-\beta \check A)^{-1} z_0 -  \bar b_0 ]  \EQN guess103 $$
so that $b(z_0, \bar c) = \bar b_0$ also satisfies equation \Ep{guess102}.
 
To begin, we note that the consumer's budget constraint is now
$$ c_t + b_{t-1}(z_t) = \int  q(z_{t+1} | z_t ) b_t(z_{t+1}) d z_{t+1} + y_t  \EQN CMbudget101 $$
where $q(z_{t+1} | z_t)$ is a pricing kernel for one-period  state contingent claims and $b_t(z_{t+1})$ is  now a vector of state-contingent debts due in state $z_{t+1}$
 that are chosen at time $t$.
To execute our reverse engineering exercise, we exogenously set the pricing kernel to 
$$ q(z_{t+1} | z_t ) = \beta \phi(z_{t+1} \ z_t) , \EQN kernel101 $$
where $\phi(z_{t+1} \ z_t)$ is the probability density of $z_{t+1}$ conditional on $z_t$.  This specification of the pricing kernel
implies that the price of a risk-free bond must be
$$ \int q(z_{t+1} | z_t) d z_{t+1} = \beta \int \phi(z_{t+1} \ z_t) = \beta $$
so that the gross interest rate on one-period risk free loans is  $\beta^{-1}$ as we want.




To construct a financial plan that supports constant consumption, temporarily take $\bar c$ as known, substitute \Ep{guess101} into
\Ep{stareatit101} and solve for a state-dependent debt function $b(z_t)$:
$$ b(z_t, \bar c) =  \check G(I-\beta \check A)^{-1} z_t - {\frac{1}{1-\beta}} \bar c . \EQN guess102 $$
To pin down the constant $\bar c$, we assume that the consumer enters time $0$ owing debt $\bar b_0$ and that $\bar c$ solves  the following version of \Ep{guess102} at time $t=0$:
$$ \bar c = (1 -\beta) [ \check G(I-\beta \check A)^{-1} z_0 -  \bar b_0 ]  \EQN guess103 $$
so that $b(z_0, \bar c) = \bar b_0$ also satisfies equation \Ep{guess102}.
 




Notice that $\bar c$ in equation \Ep{guess103} equals $c_0$ given by the optimal decision rule \Ep{eqn:lucascritiquecons} for the incomplete markets version
of the LQ consumption smoothing model.   From equation \Ep{sprob16;a}, we know that in the incomplete markets version of the model
$$ c_{t+1} = c_t + b w_{t+1} $$
where $b = (1-\beta) \check G (I - \beta \check A)^{-1} \check C $
so that 
$$ \eqalign{ E_0 c_t & = c_0 \cr
              E_0(c_t  - c_0)^2 & = t b b' } $$
Because the consumer dislikes variance in consumption, he prefers the completely smoothed   consumption stream  that he can choose in the complete markets version of the model.


\medskip


\bye
